\documentclass[10pt,a4paper]{amsart}
\usepackage[margin=19mm]{geometry}
\usepackage{amsmath,amssymb,mathtools}
\usepackage[scr=rsfs,cal=euler]{mathalfa}
\usepackage{xfrac}
\usepackage[unicode,hidelinks,bookmarks=false]{hyperref}
\newcommand{\ie}{i.e.\ }
\newcommand{\cf}{cf.\ }
\newcommand{\resp}{respectively\ }
\newcommand{\Subsection}{\subsection}
\providecommand{\nobreakdash}{\nobreak\hbox{-}}
\setlength{\emergencystretch}{2em}
\multlinegap0pt
\usepackage{setspace}
\usepackage{braket}
\usepackage{latexsym}
\usepackage{textcomp}
\usepackage{mathtools}
\usepackage{here}
\newtheorem{Thm}{Theorem}[section]
\newtheorem{Que}[Thm]{Question}
\newtheorem{Lem}[Thm]{Lemma}
\newtheorem{Prop}[Thm]{Proposition}
\newtheorem{Conj}[Thm]{Conjecture}
\newtheorem{Cor}[Thm]{Corollary}
\newtheorem{Cla}[Thm]{Claim}
\newtheorem{Fac}[Thm]{Fact}
\newtheorem{Axi}[Thm]{Axiom}
\newtheorem{Ass}[Thm]{Assumption}
\newtheorem{Def}[Thm]{Definition}
\newtheorem{Rem}[Thm]{Remark}
\newtheorem{Exa}[Thm]{Example}
\begin{document}
\title[Asymptotic estimates for multiple point ranges of transient ]{Asymptotic estimates for multiple point ranges of transient random walks on graphs}
\author{Kazuki Okamura}
\date{}
\maketitle
\noindent\emph{Source and licence.} Asymptotic estimates for multiple point ranges of transient random walks on graphs. Version arXiv:2606.00648v1 (CC BY 4.0). This material is distributed under the Creative Commons Attribution 4.0 International licence, http://creativecommons.org/licenses/by/4.0/, as recorded in the arXiv OAI licence metadata for this exact version and shown on the abstract page https://arxiv.org/abs/2606.00648v1. Full rights notice: licenses/arxiv-2606.00648-CC-BY-4.0.txt.\par\smallskip
\noindent\textit{Editorial scope.} Counted unit: the original \texttt{Thm} environment \texttt{thm:func-loc-time-Lp} at source lines 515--527 together with its complete proof at lines 529--560; the delivered document keeps source lines 128--671 so that every referenced label and every citation resolves inside it. Native editable source is the paper's own concatenated TeX; the AMS wrapper is editorial formatting only. No publisher class, upstream script or unrelated artwork is redistributed.
\begin{Thm}\label{thm:main-exp}
If $(U_0)$ holds, then, for every $j \ge 1$, 
\begin{align*} 
(F_{\inf})^{j-1} (1-F_{\sup}) &\le  \liminf_{n \to \infty} \frac{\inf_{x \in \mathcal{X}} E^{x} \left[R_n^j \right]}{n} \\
&\le \limsup_{n \to \infty} \frac{\sup_{x \in \mathcal{X}} E^{x} \left[R_n^j \right]}{n}  \le (F_{\sup})^{j-1} (1-F_{\inf}) 
\end{align*}
and 
\begin{align*}  
(F_{\inf})^{j-1} (1-F_{\sup})^2 &\le  \liminf_{n \to \infty} \frac{\inf_{x \in \mathcal{X}} E^{x} \left[R_n^{(j)} \right]}{n} \\
&\le \limsup_{n \to \infty} \frac{\sup_{x \in \mathcal{X}} E^{x} \left[R_n^{(j)} \right]}{n}  \le (F_{\sup})^{j-1} (1-F_{\inf})^2. 
\end{align*}
\end{Thm}

If  the Markov chain $(X_n)_n$ is the simple random walk on a vertex-transitive graph or the random walk on a countable group, then $F_{\sup} = F_{\inf}$. 
If  the Markov chain $(X_n)$ is recurrent, then $F_{\sup} = F_{\inf} = 1$. 
$(U_0)$ is identical with the condition in \cite[Lemma 2.1]{Okamura2014}. 

\begin{Thm}\label{thm:main-SLLN}
If $(U_1)$ holds, then, for every $j \ge 1$ and every $x \in \mathcal{X}$, 
the following statements hold $P^x$-a.s.: 
\begin{equation}\label{eq:upper-SLLN-1}
\limsup_{n \to \infty} \frac{R_n^j}{n}  \le (F_{\sup})^{j-1} (1-F_{\inf}). \ 
\end{equation}
\begin{equation}\label{eq:upper-SLLN-2}
\limsup_{n \to \infty} \frac{R_n^{(j)}}{n}  \le (F_{\sup})^{j-1} (1-F_{\inf})^2. 
\end{equation} 
\begin{equation}\label{eq:lower-SLLN-1}
\liminf_{n \to \infty} \frac{R_n^j}{n}  \ge (F_{\inf})^{j-1} (1-F_{\sup}). \ 
\end{equation}
\begin{equation}\label{eq:lower-SLLN-2}
\liminf_{n \to \infty} \frac{R_n^{(j)}}{n}  \ge (F_{\inf})^{j-1} (1-F_{\sup})^2. 
\end{equation} 
\end{Thm}

In Section \ref{sec:examples}, we give a sufficient condition for $(U_1)$ in terms of heat kernel. 
To our knowledge, the Nash-type on-diagonal heat kernel upper bounds needed to verify $(U_1)$ are not known to hold for all transient random walks on countable groups.
However, under additional assumptions, Nash-type bounds and related estimates have been studied extensively; see Saloff-Coste and Zheng  \cite{SCZ2016} for example. 
We remark that there exist infinite vertex-transitive graphs which are not Cayley graphs of groups, as shown by Woess \cite{Woess2013}. 

\section{Proofs}

In this section, we prove Theorems \ref{thm:main-exp} and \ref{thm:main-SLLN}. 
As an outline, we follow \cite{Pitt1974}; in particular, we use induction on $j$ in the estimates of expectations. 
However, we cannot apply the ergodic theorem and several parts of the argument have to be changed significantly. 
The main difficulty is to estimate $E_{m,n}$ defined in \eqref{eq:Emn} below.
In \cite{Pitt1974},  this part is handled by the ergodic theorem; here we use the uniform tail condition $(U_1)$ instead. 

\subsection{Proof of Theorem \ref{thm:main-exp}} 

Let $\mathcal{A}_{k}^{n} \coloneqq \left\{X_k \ne X_{\ell}, k+1 \le \ell \le n \right\}$ for $k < n$ and $\mathcal{A}_{n}^{n}$ be the whole event. 
Let $\mathcal{B}_{k}^{(j)} \coloneqq \left\{   |\{i \in \{1,\dots,k\} : X_i = X_{k} \}| =  j \right\}$. 
We remark that $\mathcal{B}_{k}^{(j)} = \emptyset$ if $j > k$.  
Then, 
\[ R^{(j)}_{n} = \sum_{k=1}^{n} {\bf 1}_{\mathcal{B}_{k}^{(j)} \cap \mathcal{A}_{k}^{n}}, \textup{ and } R^{j}_{n} = \sum_{k=1}^{n} {\bf 1}_{\mathcal{B}_{k}^{(j)}}. \] 
Hence, for every $x \in \mathcal{X}$, 
\begin{equation}\label{eq:basic-exp-1} 
E^x \left[ R^{(j)}_{n} \right] = \sum_{k=1}^{n} P^x \left(\mathcal{B}_{k}^{(j)} \cap \mathcal{A}_{k}^{n}\right) \textup{ and } E^x \left[ R^{j}_{n} \right] = \sum_{k=1}^{n} P^x \left(\mathcal{B}_{k}^{(j)} \right). 
\end{equation}

For $n \ge 0$, let  
\[ F_{\inf} (n) \coloneqq \inf_{x \in \mathcal{X}} P^x (T_x \le n) \ \textup{ and } \   F_{\sup} (n) \coloneqq \sup_{x \in \mathcal{X}} P^x (T_x \le n).  \]

Then, by the Markov property, 
\begin{equation}\label{eq:basic-Markov-ineq} 
1 - F_{\sup}(n-k) \le P^x \left( \mathcal{A}_{k}^{n} \, \middle| \, \mathcal{B}_{k}^{(j)} \right) \le 1 - F_{\inf}(n-k), \ x \in \mathcal{X}. 
\end{equation}


First, we deal with the upper bounds. 
By \eqref{eq:basic-exp-1} and \eqref{eq:basic-Markov-ineq}, 
\[ E^x \left[ R^{(j)}_{n} \right] \ge \sum_{k=1}^{n} P^x \left(\mathcal{B}_{k}^{(j)} \right) (1- F_{\sup}(n-k)), \ x \in \mathcal{X}. \]
Using this, \eqref{eq:basic-exp-1}, and the fact that $F_{\sup}(m) \le F_{\sup}$ for each $m \ge 1$,   
it holds that 
\begin{equation}\label{eq:comparision-ineq-1} 
E^x \left[ R^{(j)}_{n} \right] \ge (1-F_{\sup}) E^x \left[ R^{j}_{n} \right], \ x \in \mathcal{X}. 
\end{equation}
Since $R_n^j - R_n^{j+1} = R_n^{(j)}$, 
\[ E^x \left[ R^{j+1}_{n} \right] \le F_{\sup} E^x \left[ R^{j}_{n} \right], \ x \in \mathcal{X}, \]
and hence, 
\begin{equation}\label{eq:exp-unif-upperbound} 
\sup_{x \in \mathcal{X}} E^x \left[ R^{j+1}_{n} \right] \le F_{\sup}  \sup_{x \in \mathcal{X}} E^x \left[ R^{j}_{n} \right]. 
\end{equation}

We see that 
\[ \limsup_{n \to \infty} \frac{\sup_{x \in \mathcal{X}} E^{x} \left[R_n \right]}{n} \le 1 - F_{\inf}. \]
Hence, by induction on $j$, 
\begin{equation}\label{eq:upper-exp-1} 
\limsup_{n \to \infty} \frac{\sup_{x \in \mathcal{X}} E^{x} \left[R_n^j \right]}{n} \le (F_{\sup})^{j-1} (1 - F_{\inf}). 
\end{equation}

We consider $E^x \left[ R^{(j)}_{n} \right]$. 
By \eqref{eq:basic-exp-1} and \eqref{eq:basic-Markov-ineq}, 
\begin{equation}\label{eq:upper-exp-2}  
E^x \left[ R^{(j)}_{n} \right] \le  \sum_{k=1}^{n} P^x \left(\mathcal{B}_{k}^{(j)} \right) (1- F_{\inf}(n-k)),  \ x \in \mathcal{X}.  
\end{equation}
By $(U_0)$, it holds that 
\[ \lim_{n \to \infty} F_{\inf} (n) = F_{\inf} \ \textup{ and } \  \lim_{n \to \infty} F_{\sup} (n) = F_{\sup}. \]

Let $\epsilon > 0$. 
Let $N_{\epsilon}$ be an integer such that for every $n > N_{\epsilon}$, $F_{\inf}(n) \ge F_{\inf} - \epsilon$. 
Then, for every $n > N_{\epsilon}$ and $x \in \mathcal{X}$, 
\begin{equation*}
\sum_{k=1}^{n} P^x \left(\mathcal{B}_{k}^{(j)} \right) (1- F_{\inf}(n-k)) \le (1 - F_{\inf} + \epsilon) \sum_{k=1}^{n - N_{\epsilon}}  P^x \left(\mathcal{B}_{k}^{(j)} \right) + N_{\epsilon}. 
\end{equation*}
Combining this with \eqref{eq:upper-exp-2} and \eqref{eq:basic-exp-1}, 
\begin{equation}\label{eq:comparision-ineq-2-1} 
E^x \left[ R^{(j)}_{n} \right]  \le  (1 - F_{\inf} + \epsilon) E^x \left[ R^{j}_{n} \right]  + N_{\epsilon}, \ n > N_{\epsilon}, \ x \in \mathcal{X}.  
\end{equation}
By this and \eqref{eq:upper-exp-1},
\[ \limsup_{n \to \infty} \frac{\sup_{x \in \mathcal{X}} E^{x} \left[R_n^{(j)} \right]}{n}  \le (1 - F_{\inf} + \epsilon)  (F_{\sup})^{j-1} (1 - F_{\inf}).   \]
Letting $\epsilon \to +0$, 
\[ \limsup_{n \to \infty} \frac{\sup_{x \in \mathcal{X}} E^{x} \left[R_n^{(j)} \right]}{n}  \le (F_{\sup})^{j-1} (1 - F_{\inf})^2.   \]

We now deal with the lower bound. 
By \eqref{eq:comparision-ineq-2-1} and $R_n^j - R_n^{j+1} = R_n^{(j)}$, 
\[ (F_{\inf} - \epsilon) E^x \left[ R^{j}_{n} \right] \le E^x \left[ R^{j+1}_{n} \right] + N_{\epsilon}, \ n > N_{\epsilon}, \ x \in \mathcal{X}. \]
Hence, 
\[ (F_{\inf} - \epsilon) \inf_{x \in \mathcal{X}} E^x \left[ R^{j}_{n} \right] \le \inf_{x \in \mathcal{X}} E^x \left[ R^{j+1}_{n} \right] + N_{\epsilon}, \ n > N_{\epsilon}. \]
Dividing by $n$ and letting $n \to \infty$, 
\[ (F_{\inf} - \epsilon) \liminf_{n \to \infty} \frac{\inf_{x \in \mathcal{X}} E^{x} \left[R_n^j \right]}{n} \le \liminf_{n \to \infty} \frac{\inf_{x \in \mathcal{X}} E^{x} \left[R_n^{j+1} \right]}{n}. \]
Letting $\epsilon \to +0$, 
\[ F_{\inf} \liminf_{n \to \infty} \frac{\inf_{x \in \mathcal{X}} E^{x} \left[R_n^j \right]}{n} \le \liminf_{n \to \infty} \frac{\inf_{x \in \mathcal{X}} E^{x} \left[R_n^{j+1} \right]}{n}. \]

By using the last-exit decomposition as in the proof of  \cite[Theorem 1.2]{Okamura2014}, 
we obtain that
\begin{equation}\label{eq:exp-lower-most-basic} 
E^{x} \left[R_n \right] \ge \sum_{k=1}^{n} \inf_{y \in \mathcal{X}} P^y (T_y \ge k) \ge n \inf_{y \in \mathcal{X}} P^y (T_y = \infty) = n(1 - F_{\sup}). 
\end{equation}

It follows that 
\[ \liminf_{n \to \infty} \frac{\inf_{x \in \mathcal{X}} E^{x} \left[R_n \right]}{n} \ge 1 - F_{\sup}. \]
Therefore, by induction on $j$, 
\begin{equation*}\label{eq:lower-exp-1} 
\liminf_{n \to \infty} \frac{\inf_{x \in \mathcal{X}} E^{x} \left[R_n^j \right]}{n} \ge (F_{\inf})^{j-1} (1 - F_{\sup}). 
\end{equation*}
By this and \eqref{eq:comparision-ineq-1}, 
\[ \liminf_{n \to \infty} \frac{\inf_{x \in \mathcal{X}} E^{x} \left[R_n^{(j)} \right]}{n} \ge (F_{\inf})^{j-1} (1 - F_{\sup})^2. \]
This completes the proof. 

\subsection{Proof of Theorem \ref{thm:main-SLLN}}


We first deal with the upper estimates. 

Let 
\[ T^{j}_{k,n} \coloneqq \left|\left\{x \in \mathcal{X} \colon \ell(kn,x) - \ell((k-1)n,x) \ge  j \right\}\right|  \]
and 
\[ T^{(j)}_{k,n} \coloneqq \left|\left\{x \in \mathcal{X} \colon  \ell(kn,x) - \ell((k-1)n,x) =  j \right\}\right|.  \]
Let 
$\mathcal{R}_{k,n} \coloneqq \left\{X_{i} \colon (k-1)n + 1 \le i \le kn \right\}$  
and  
\begin{equation}\label{eq:Emn} 
E_{m,n} \coloneqq \left|\bigcup_{k_1, k_2 \in \{1, \dots, m\}, k_1 \ne k_2} \mathcal{R}_{k_1, n} \cap \mathcal{R}_{k_2, n} \right|. 
\end{equation}

Then, 
\[ R_{mn}^{j} \le \sum_{k=1}^{m}   T^{j}_{k,n} + E_{m,n}, \ \textup{ and }  \  R_{mn}^{(j)} \le \sum_{k=1}^{m}   T^{(j)}_{k,n} + E_{m,n}. \]
For notational convenience, let 
\[ F^{(u)}_{j} \coloneqq (F_{\sup})^{j-1} (1-F_{\inf}), \ \textup{ and }  \ F^{(u)}_{(j)} \coloneqq (F_{\sup})^{j-1} (1-F_{\inf})^2. \]
Then, for every $\epsilon > 0$, 
\[ P^x \left(R_{mn}^{j} \ge mn(F_j^{(u)} + \epsilon)\right) \le P^x \left(  \sum_{k=1}^{m}   T^{j}_{k,n} \ge mn \left(F_j^{(u)} + \frac{\epsilon}{2} \right) \right) + P^x \left( E_{m,n} \ge mn \frac{\epsilon}{2} \right). \]

We first give an upper bound for  $\displaystyle P^x \left(  \sum_{k=1}^{m}   T^{j}_{k,n} \ge mn \left(F_j^{(u)} + \frac{\epsilon}{2} \right) \right)$. 

\begin{Lem}\label{lem:Laplace-estimate-bdd}
For every $\epsilon > 0$, there exists $\theta(\epsilon) > 0$ such that 
for every random variable $Y$ on a probability space satisfying $|Y| \le 1$, 
\begin{equation}\label{eq:every-bdd-rv-Laplace}
E\left[\exp(\theta(\epsilon) Y) \right] \le 1 + \theta(\epsilon) (E[Y] + \epsilon) \le \exp\left(  \theta(\epsilon) (E[Y] + \epsilon) \right). 
\end{equation}
\end{Lem}

\begin{proof}
By the Lebesgue convergence theorem and the fact that $E[Y^n] \le 1$, 
\[  E\left[\exp(\theta(\epsilon) Y) \right] = \sum_{n=0}^{\infty} \frac{\theta^n}{n!} E[Y^n] \le 1 + \theta E[Y] + \frac{\theta^2}{2} \exp(\theta).    \]
Let $\theta(\epsilon)$ be a positive constant such that 
$\theta \exp(\theta) = \epsilon$. 
Then, using the inequality $1+x \le \exp(x)$, 
we have \eqref{eq:every-bdd-rv-Laplace}. 
\end{proof}


\begin{Lem}\label{lem:Laplace-estimate-multi-range}
For every $\epsilon > 0$, there exists  $N(\epsilon) \in \mathbb{N}$ such that 
for every $n \ge N(\epsilon)$, 
\begin{equation}\label{eq:multi-range-Laplace}
\sup_{x \in \mathcal{X}} E^x \left[\exp\left(\frac{\theta(\epsilon/2)}{n} R_n^j \right) \right] \le \exp\left(\theta(\epsilon/2) (F_j^{(u)} + \epsilon) \right),
\end{equation}
where $\theta (\cdot)$ is the function appearing in Lemma \ref{lem:Laplace-estimate-bdd}.  
\end{Lem}

\begin{proof}
Since $0 \le R^j_n \le n$, 
by Lemma \ref{lem:Laplace-estimate-bdd}, 
\[ \sup_{x \in \mathcal{X}} E^x \left[\exp\left(\frac{\theta(\epsilon/2)}{n} R_n^j \right) \right] \le \exp\left(  \theta(\epsilon/2) \left(\frac{E^x \left[R_n^j \right] }{n} + \frac{\epsilon}{2} \right) \right). \]
By Theorem \ref{thm:main-exp}, 
there exists $N(\epsilon) \in \mathbb{N}$ such that 
for every $n \ge N(\epsilon)$, 
$\displaystyle \frac{E^x \left[R_n^j \right] }{n} \le F_j^{(u)} + \frac{\epsilon}{2}$. 
Hence, 
\eqref{eq:multi-range-Laplace} holds for every $n \ge N(\epsilon)$. 
\end{proof}


\begin{Lem}\label{lem:upper-Markov-tail}
For every $\epsilon > 0$, there exist $C(\epsilon) > 0$  and $N_1 (\epsilon) \in \mathbb{N}$ such that 
for every $n \ge N_1 (\epsilon)$ and every $m \in \mathbb{N}$, 
\begin{equation*}\label{eq:multi-range-tail}
\sup_{x \in \mathcal{X}} P^x \left(  \sum_{k=1}^{m}   T^{j}_{k,n} \ge mn \left(F_j^{(u)} + \frac{\epsilon}{2} \right) \right)  \le \exp\left(-C(\epsilon) m \right). 
\end{equation*}
\end{Lem}

\begin{proof}
By the Markov property, for every $\lambda > 0$ and $x \in \mathcal{X}$, 
\[  P^x \left(  \sum_{k=1}^{m}   T^{j}_{k,n} \ge mn \left(F_j^{(u)} + \frac{\epsilon}{2} \right) \right) \le \exp\left(-\lambda mn \left(F_j^{(u)} + \frac{\epsilon}{2} \right) \right) E^x \left[ \prod_{k=1}^{m} \exp\left( \lambda T^{j}_{k,n}  \right)  \right] \]
\[ \le \left( \exp\left(-\lambda n \left(F_j^{(u)} + \frac{\epsilon}{2} \right) \right) \sup_{x \in \mathcal{X}} E^x \left[ \exp(\lambda R^j_n) \right] \right)^m. \]
Let $n \ge N(\epsilon/4)$ and $\lambda = \frac{1}{n} \theta(\epsilon/8)$. 
Then, by Lemma \ref{lem:Laplace-estimate-multi-range}, 
\[  \exp\left(-\lambda n \left(F_j^{(u)} + \frac{\epsilon}{2} \right) \right) \sup_{x \in \mathcal{X}} E^x \left[ \exp(\lambda R^j_n) \right] \le \exp\left(- \frac{\epsilon}{4} \theta(\epsilon/8)\right). \]
Thus the assertion holds for $\displaystyle C(\epsilon) \coloneqq \frac{\epsilon}{4} \theta(\epsilon/8)$ and $N_1 (\epsilon) \coloneqq N(\epsilon/4)$.
\end{proof}

Second, we give an upper bound for  $\displaystyle P^x \left( E_{m,n} \ge mn \frac{\epsilon}{2} \right)$. 
Let 
\[ \mathcal{A}_{k}^{\infty} \coloneqq \left\{X_{\ell} \ne X_k,  \ell \ge k+1 \right\} = \bigcap_{n \ge k+1} \mathcal{A}_{k}^n.\]  
Then, by the same argument as in \cite[Proof of Lemma]{Pitt1974}, 
it holds that for $n^{\prime} < n$, 
\[ E_{m,n} \le \sum_{k=1}^{mn} {\bf 1}_{\mathcal{A}_{k}^{k+n^{\prime}} \setminus \mathcal{A}_{k}^{\infty}} + mn^{\prime}. \]

Henceforth, we denote the integer part of a real number $x$ by $\lfloor x \rfloor$. 
For $n \ge 2$, let 
\begin{equation}\label{eq:def-n-prime}
n^{\prime} \coloneqq \left\lfloor \frac{n}{\log n} \right\rfloor. 
\end{equation}
Then, there exists $N_2 (\epsilon) \in \mathbb{N}$ such that for every $n \ge N_2 (\epsilon)$, $n^{\prime}/n < \epsilon/4$. 
Hence, it holds that for every $n \ge N_2 (\epsilon)$ and every $m \ge 1$, 
\begin{equation}\label{eq:Emn-upper} 
P^x \left( E_{m,n} \ge mn \frac{\epsilon}{2} \right) \le P^x \left(  \sum_{k=1}^{mn} {\bf 1}_{\mathcal{A}_{k}^{k+n^{\prime}} \setminus \mathcal{A}_{k}^{\infty}} \ge mn \frac{\epsilon}{4} \right) \le \frac{4}{\epsilon} \sup_{x \in \mathcal{X}} P^x (n^{\prime} < T_x < \infty). 
\end{equation}

By $(U_1)$, there exists a constant $C_0$ such that 
\begin{equation}\label{eq:c0}
\sup_{x \in \mathcal{X}} P^x (n < T_x < \infty) \le C_0 (\log n)^{-1-\delta}, \ \ n \ge 1.  
\end{equation}

Thus, 
for every $\epsilon > 0$, every $n \ge N_1 (\epsilon) + N_2 (\epsilon)$ and every $m \in \mathbb{N}$, 
it follows that 
\begin{equation*}
P^x \left(R_{mn}^{j} \ge mn(F_j^{(u)} + \epsilon)\right) \le \exp\left(-C(\epsilon) m \right) + \frac{4C_0}{\epsilon} (\log n - \log \log n)^{-1-\delta}. 
\end{equation*}

Let $\eta > 0$. 



For $m_k = r_k =  \lfloor (1+\eta)^{k/2} \rfloor$, 
\[ \exp\left(-C(\epsilon) m_k \right) + \frac{4C_0}{\epsilon} (\log r_k - \log \log r_k)^{-1-\delta} = O\left(k^{-1-\delta}\right), \ k \to \infty. \]
Let $a_k \coloneqq m_k r_k$. 
Then, 
 \[ \sum_{k=1}^{\infty} P^x \left(R_{a_k}^{j} \ge a_k (F_j^{(u)} + \epsilon)\right) < +\infty, \]
and by the Borel-Cantelli lemma, 
\begin{equation}\label{eq:lacunary-upper} 
\limsup_{k \to \infty} \frac{R_{a_k}^{j}}{a_k} \le F_j^{(u)} + \epsilon, \ \textup{ $P^x$-a.s.}  
\end{equation}
For $a_k \le n \le a_{k+1}$, 
$\displaystyle \frac{R^j_n}{n} \le \frac{a_{k+1}}{a_k} \frac{R_{a_{k+1}}^{j}}{a_{k+1}}$. 
By this estimate and the equality $\displaystyle \lim_{k \to \infty} \frac{a_{k+1}}{a_k} = 1+\eta$, 
\[ \limsup_{n \to \infty} \frac{R_{n}^{j}}{n} \le (1+\eta)(F_j^{(u)} + \epsilon), \ \textup{ $P^x$-a.s.}  \]
By letting  $\eta \to 0$ and $\epsilon \to 0$, 
we see that  \eqref{eq:upper-SLLN-1} holds $P^x$-a.s.  


We now deal with $R^{(j)}_n$. 
In the same manner as in the derivation of \eqref{eq:lacunary-upper}, we obtain that   
\[ \limsup_{k \to \infty} \frac{R_{a_k}^{(j)}}{a_k} \le F_{(j)}^{(u)} + \epsilon, \ \textup{ $P^x$-a.s.}  \]
Since 
\[ R_n^{(j)} \le R_{a_{k+1}}^{(j)} + a_{k+1} - a_k,  \ a_k \le n \le a_{k+1},\]  
it holds that 
\[ \limsup_{n \to \infty} \frac{R_n^{(j)}}{n}  \le \limsup_{k \to \infty} \frac{a_{k+1}}{a_k} \frac{R_{a_{k+1}}^{(j)}}{a_{k+1}} + \frac{a_{k+1} - a_k}{a_k} \le (1+\eta)(F_{(j)}^{(u)} + \epsilon) + \eta, \ \textup{ $P^x$-a.s.}   \]
By letting  $\eta \to 0$ and $\epsilon \to 0$, 
we see that  \eqref{eq:upper-SLLN-2} holds $P^x$-a.s.  


We now deal with the lower estimates. 

We see that 
\[ R_{mn}^{j} \ge \sum_{k=1}^{m}  (T^{j}_{k,n} - E_{m,n}), \ \textup{ and }  \  R_{mn}^{(j)} \ge \sum_{k=1}^{m} (T^{(j)}_{k,n} - E_{m,n}). \]
For ease of notation, let 
\[ F^{(l)}_{j} \coloneqq (F_{\inf})^{j-1} (1-F_{\sup}), \ \textup{ and }  \ F^{(l)}_{(j)} \coloneqq (F_{\inf})^{j-1} (1-F_{\sup})^2. \]
Then, for every $\epsilon > 0$, 
\[ P^x \left(R_{mn}^{j} \le mn(F_j^{(l)} - \epsilon)\right) \le P^x \left(  \sum_{k=1}^{m}   T^{j}_{k,n} \le mn \left(F_j^{(l)} - \frac{\epsilon}{2} \right) \right) + P^x \left(  E_{m,n} \ge n \frac{\epsilon}{2} \right). \]


We first give an upper bound for  $\displaystyle P^x \left(  \sum_{k=1}^{m}   T^{j}_{k,n} \le mn \left(F_j^{(l)} - \frac{\epsilon}{2} \right) \right)$. 
By Theorem \ref{thm:main-exp}, 
it holds that 
for every $\epsilon > 0$, 
there exists $N(\epsilon) \in \mathbb{N}$ such that 
for every $n \ge N(\epsilon)$, 
$\displaystyle \frac{E^x \left[R_n^j \right] }{n} \ge F_j^{(l)} - \frac{\epsilon}{4}$, which is equivalent with $\displaystyle \frac{E^x \left[n - R_n^j \right] }{n} \le 1 - F_j^{(l)} + \frac{\epsilon}{4}$. 

In the same manner as in the proof of Lemma \ref{lem:Laplace-estimate-multi-range}, 
we can show that 
for every $\epsilon > 0$, there exists  $N(\epsilon) \in \mathbb{N}$ such that 
for every $n \ge N(\epsilon)$, 
\[ \sup_{x \in \mathcal{X}} E^x \left[\exp\left(\frac{\theta(\epsilon/8)}{n} (n-R_n^j) \right) \right] \le \exp\left(\theta(\epsilon/8) \left(1 - F_j^{(l)} + \frac{3\epsilon}{8} \right) \right).\] 


Since $T^{j}_{k,n} \le n$, 
by the Markov property, for every $\lambda > 0$ and $x \in \mathcal{X}$, 
\[  P^x \left(  \sum_{k=1}^{m}   T^{j}_{k,n} \le mn \left(F_j^{(l)} - \frac{\epsilon}{2} \right) \right) = P^x \left(  \sum_{k=1}^{m}   (n-T^{j}_{k,n}) \ge mn \left(1 - F_j^{(l)} + \frac{\epsilon}{2} \right) \right)\]
\[ \le \exp\left(-\lambda mn \left(1 - F_j^{(l)} + \frac{\epsilon}{2} \right) \right) E^x \left[ \prod_{k=1}^{m} \exp\left( \lambda (n-T^{j}_{k,n})  \right)  \right] \]
\[ \le \left( \exp\left(-\lambda n \left(1 - F_j^{(l)} + \frac{\epsilon}{2} \right) \right) \sup_{x \in \mathcal{X}} E^x \left[ \exp\left(\lambda (n-R^j_n) \right) \right] \right)^m. \]

In the same manner as in the proof of Lemma \ref{lem:upper-Markov-tail}, 
we can show that 
for every $\epsilon > 0$, there exist $C(\epsilon) > 0$  and $N_3 (\epsilon) \in \mathbb{N}$ such that 
for every $n \ge N_3 (\epsilon)$ and every $m \in \mathbb{N}$, 
\begin{equation*}
\sup_{x \in \mathcal{X}} P^x \left(  \sum_{k=1}^{m}   T^{j}_{k,n} \le mn \left(F_j^{(l)} - \frac{\epsilon}{2} \right) \right)  \le \exp\left(-C(\epsilon) m \right). 
\end{equation*}

Second, we give an upper bound for $\displaystyle P^x \left( E_{m,n} \ge n \frac{\epsilon}{2} \right)$. 
Let $n^{\prime}$ be as in \eqref{eq:def-n-prime}. 
As in \eqref{eq:Emn-upper}, by using $(U_1)$,  
we see that there exists $N_4 (\epsilon) \in \mathbb{N}$ such that for every $n \ge N_4 (\epsilon)$ and $m \in \mathbb{N}$ satisfying $mn^{\prime}/n < \epsilon/4$, 
\begin{equation*}\label{eq:Emn-upper-2} 
P^x \left( E_{m,n} \ge n \frac{\epsilon}{2} \right) \le \frac{2m}{\epsilon - 2m n^{\prime}/n}  \sup_{x \in \mathcal{X}} P^x (n^{\prime} < T_x < \infty). 
\end{equation*}



Recall \eqref{eq:c0}. 
Thus we obtain that 
for every $\epsilon > 0$, 
for every $n \ge N_3 (\epsilon) + N_4 (\epsilon)$ and every $m \in \mathbb{N}$ satisfying $mn^{\prime}/n < \epsilon/4$, 
\begin{equation*}
P^x \left(R_{mn}^{j} \le mn(F_j^{(l)} - \epsilon)\right) \le \exp\left(-C(\epsilon) m \right) + \frac{4C_0}{\epsilon} m (\log n - \log \log n)^{-1-\delta}. 
\end{equation*}

Let $\eta > 0$. 
For $m_k = \lfloor (\log k)^2 \rfloor$ and $r_k = \lfloor (1+\eta)^k \rfloor$, 
it holds that $m_k (r_k)^{\prime} / r_k < \epsilon/4$ for large $k$, and, 
\[ \exp\left(-C(\epsilon) m_k \right) + \frac{4C_0}{\epsilon} m_k (\log r_k - \log \log r_k)^{-1-\delta} = O\left(k^{-1- \frac{\delta}{2}} \right), \ k \to \infty. \] 

Let $b_k \coloneqq m_k r_k$. 
Then, 
\[ \sum_{k=1}^{\infty} P^x \left(R_{b_k}^{j} \le b_k (F_j^{(l)} - \epsilon)\right) < +\infty. \]
Hence, by the Borel-Cantelli lemma, 
\[ \liminf_{k \to \infty} \frac{R_{b_k}^{j}}{b_k} \ge F_j^{(l)} - \epsilon, \ \textup{ $P^x$-a.s.}  \]
It holds that for $b_k \le n \le b_{k+1}$, 
$\displaystyle \frac{R^j_n}{n} \ge \frac{b_{k}}{b_{k+1}} \frac{R_{b_{k}}^{j}}{b_{k}}$. 
By this estimate and the equality $\displaystyle \lim_{k \to \infty} \frac{b_{k}}{b_{k+1}} = \frac{1}{1+\eta}$, 
\begin{equation}\label{eq:lacunary-lower} 
\liminf_{n \to \infty} \frac{R_{n}^{j}}{n} \ge \frac{F_j^{(l)} - \epsilon}{1+\eta}, \ \textup{ $P^x$-a.s.}  
\end{equation}
Letting  $\eta \to 0$ and $\epsilon \to 0$, 
we see that  \eqref{eq:lower-SLLN-1} holds $P^x$-a.s.  

We now deal with $R^{(j)}_n$. 
In the same manner as in the derivation of \eqref{eq:lacunary-lower},  
\[ \liminf_{k \to \infty} \frac{R_{b_k}^{(j)}}{b_k} \ge F_{(j)}^{(l)} - \epsilon, \ \textup{ $P^x$-a.s.}  \]
Since 
\[ R_n^{(j)} \ge R_{b_{k}}^{(j)} - (b_{k+1} - b_k),  \ b_k \le n \le b_{k+1},\]  
it holds that 
\[ \liminf_{n \to \infty} \frac{R_n^{(j)}}{n}  \ge \liminf_{k \to \infty} \frac{b_{k}}{b_{k+1}} \frac{R_{b_{k}}^{(j)}}{b_{k}} - \frac{b_{k+1} - b_k}{b_k} \ge \frac{F_{(j)}^{(l)} - \epsilon}{1+\eta} - \eta, \ \textup{ $P^x$-a.s.}   \]
Letting  $\eta \to 0$ and $\epsilon \to 0$, 
we see that \eqref{eq:lower-SLLN-2} holds $P^x$-a.s.  


\section{Functions of the local times}\label{sec:local-times}

Let $f$ be a non-negative function on $\mathbb{N} \cup \{0\}$ such that $f(0) = 0$. 
Let 
\[ G_n (f) \coloneqq \sum_{x \in \mathcal{X}} f(\ell(n,x)). \]
Let 
\[ C_{\inf}(f) \coloneqq \sum_{j \ge 1} f(j) (F_{\inf})^{j-1} (1-F_{\sup})^2 \]
and 
\[ C_{\sup}(f) \coloneqq \sum_{j \ge 1} f(j) (F_{\sup})^{j-1} (1-F_{\inf})^2. \]

Since $f(0) = 0$, we see that 
\[ G_n (f) = \sum_{j \ge 1} f(j) R_n^{(j)}. \] 


The following theorem extends  \cite[Theorem 2]{CCMP2026} and the $L^2$-convergence part of \cite[Theorem 1]{AK2020}. 


\begin{Thm}\label{thm:func-loc-time-Lp}
Let $p \ge 1$. 
Assume that  $(U_1)$ holds and 
\[ \sum_{j \ge 1} \frac{f(j)^p}{j^{p-1}} (F_{\sup})^{j} < \infty.  \]
Then 
\begin{equation}\label{eq:L2-sup-upper} 
\lim_{n \to \infty} E^x \left[ \left(\left( \frac{G_n (f)}{n} - C_{\sup}(f) \right)_{+}\right)^p \right] = 0 
\end{equation} 
and 
\begin{equation}\label{eq:L2-inf-lower}   
\lim_{n \to \infty} E^x \left[ \left(\left( \frac{G_n (f)}{n} - C_{\inf}(f) \right)_{-}\right)^p \right] = 0. 
\end{equation} 
\end{Thm}

\begin{proof}
We remark that for every $a, b \in \mathbb{R}$, 
\[ ((a+b)_{+})^p \le (a_{+} + b_{+})^p \le 2^{p-1}((a_{+})^p + (b_{+})^p). \]

Let $f_N (j) \coloneqq f(j) {\bf 1}_{\{1,\dots, N\}}(j)$. 
Then, using $C_{\sup}(f) \ge C_{\sup}(f_N)$, 
\[ \left( \frac{G_n (f)}{n} - C_{\sup}(f) \right)_+ \le \left( \frac{G_n (f_N)}{n} - C_{\sup}(f) \right)_{+} + \frac{G_n (f) - G_n (f_N)}{n}.  \]

By \eqref{eq:exp-unif-upperbound}, 
\begin{equation}\label{eq:exp-upper-unif} 
\sup_x E^x \left[R_n^{(j)} \right] \le \sup_x E^x \left[R_n^{j} \right] \le (F_{\sup})^{j-1} \sup_x E^x \left[R_n \right] \le n (F_{\sup})^{j-1}. 
\end{equation}

Let $q$ be  the conjugate exponent of $p$, i.e., $\frac{1}{p} + \frac{1}{q} = 1$, with the convention $q=\infty$ if $p=1$. 
Using the H\"older inequality  and $n = \sum_{j \ge 1} j R_n^{(j)}$, 
\[ \sum_{j \ge N+1} f(j) R_n^{(j)} = \sum_{j \ge N+1} \frac{f(j)}{j} (j R_n^{(j)})^{1/p} \cdot (j R_n^{(j)})^{1/q} \le n^{1/q} \left(\sum_{j \ge N+1} \frac{f(j)^p}{j^{p-1}} R_n^{(j)}\right)^{1/p}. \]
Hence, 
\[ \left(  \frac{G_n (f)}{n} -  \frac{G_n (f_N)}{n} \right)^p  \le \sum_{j \ge N+1} \frac{f(j)^p}{j^{p-1}} \frac{R_n^{(j)}}{n}.  \]
Using this and \eqref{eq:exp-upper-unif}, 
\[ E^x \left[ \left(  \frac{G_n (f)}{n} -  \frac{G_n (f_N)}{n} \right)^p \right] \le \sum_{j \ge N+1} \frac{f(j)^p}{j^{p-1}}(F_{\sup})^{j-1}.  \]

Since $ \frac{G_n (f_N)}{n} \le \max_{1 \le j \le N} f(j)$, 
\eqref{eq:upper-SLLN-2} and the Lebesgue convergence theorem yield that 
\[ \lim_{n \to \infty} E^x \left[ \left( \left( \frac{G_n (f_N)}{n} -  C_{\sup}(f_N)  \right)_{+} \right)^p \right] = 0. \]
Using this and $  C_{\sup}(f_N) \le   C_{\sup}(f)$, we obtain that 
\[ \lim_{n \to \infty} E^x \left[ \left( \left( \frac{G_n (f_N)}{n} -  C_{\sup}(f)  \right)_{+} \right)^p \right] = 0. \]

Hence, 
\[  \limsup_{n \to \infty} E^x \left[ \left(\left( \frac{G_n (f)}{n} - C_{\sup}(f) \right)_{+}\right)^p \right] \le 2^{p-1} \sum_{j \ge N+1} \frac{f(j)^p}{j^{p-1}}(F_{\sup})^{j-1}. \]
Letting $N \to \infty$, we have \eqref{eq:L2-sup-upper}. 
The same argument gives  \eqref{eq:L2-inf-lower}.  
\end{proof}


\begin{Prop}
Assume that $(U_1)$ holds and $F_{\inf} < 1$. 
Then for every $x \in \mathcal{X}$, 
\begin{equation}\label{eq:inf-sup-sandwich} 
C_{\inf}(f) \le \liminf_{n \to \infty} \frac{G_n (f)}{n} \le C_{\sup}(f), \ \textup{ $P^x$-a.s.}  
\end{equation}
\end{Prop}

\begin{proof}
The lower bound follows from the non-negativity of $f$ and Theorem \ref{thm:main-SLLN}. 
If $C_{\sup}(f) = \infty$, the upper bound is obvious. 
Assume that $C_{\sup}(f) < \infty$. 
By Theorem \ref{thm:func-loc-time-Lp} for $p=1$, 
there exists a sequence $(n_k)_k$ such that 
\[ \lim_{k \to \infty}  \left( \frac{G_{n_k} (f)}{n_k} - C_{\sup}(f) \right)_{+} = 0, \ \textup{ $P^x$-a.s.}  \]
This is equivalent to  
\[ \limsup_{k \to \infty}  \frac{G_{n_k} (f)}{n_k} \le C_{\sup}(f), \ \textup{ $P^x$-a.s.} \]
This implies the upper bound. 
\end{proof}

If $f(j) = j$ and $F_{\inf} = 1$, then $G_n (f) = n$ and $C_{\sup}(f) = 0$ so that \eqref{eq:inf-sup-sandwich} fails. 


The following lower bound on the expectation plays a role similar to that of  \cite[Lemma 3]{CCMP2026}. 

\begin{Lem}\label{lem:exp-lower-quantitative}
Let $x \in \mathcal{X}, n \ge 1$. 
Then, \\
(i) $E^x \left[ R_n^{(1)} \right] \ge (1-F_{\sup})^2 n$.\\
(ii) If $j \ge 2$, then for every  $\ell \in \{1,\dots, n\}$, 
\[ E^x \left[ R_n^{(j)} \right] \ge (1-F_{\sup})^2 \ell F_{\inf}\left( \left\lfloor \frac{n-\ell}{j-1} \right\rfloor\right)^{j-1}. \]
\end{Lem}

\begin{proof}
For $y \in \mathcal{X}$, 
let $ T_y^{(1)} \coloneqq T_y$ and 
\[ T_y^{(j)} \coloneqq \inf \left\{n > T_y^{(j-1)} \colon X_n = y \right\}, \ j \ge 2. \]

By the Markov property, we obtain that 
\begin{equation}\label{eq:multi-ineq-Markov}
P^y \left( T_y^{(j)} \le n \right) \ge P^y \left( T_y \le \left\lfloor \frac{n}{j} \right\rfloor \right)^{j}, \ j \ge 1, n \ge 1, y \in \mathcal{X}. 
\end{equation}
We also see that 
\begin{equation*}
E^x \left[ R_n^j \right] = \sum_{y \in \mathcal{X}}  P^x \left( T_y^{(j)} \le n \right). 
\end{equation*}
By this and the Markov property, 
\[ E^x \left[ R_n^j \right] = \sum_{y \in \mathcal{X}}  \sum_{k = 1}^{n}  P^x \left( T_y^{(j)} \le n, T_y = k \right) = \sum_{y \in \mathcal{X}}  \sum_{k = 1}^{n}  P^x \left( T_y = k \right) P^y \left( T_y^{(j-1)} \le n - k \right). \]
By this and \eqref{eq:multi-ineq-Markov}, 
\[  E^x \left[ R_n^j \right] \ge \sum_{y \in \mathcal{X}} P^x \left( T_y \le \ell \right) P^y \left( T_y^{(j-1)} \le n - \ell \right) \ge E^x \left[R_{\ell}\right] F_{\inf}\left( \left\lfloor \frac{n-\ell}{j-1} \right\rfloor\right)^{j-1}. \]
Applying \eqref{eq:comparision-ineq-1} and  \eqref{eq:exp-lower-most-basic}, we obtain the assertion. 
\end{proof}

We now consider the following stronger uniform tail  condition $(U_2)$: 
\[ \sup_{x} P^x (n < T_x < \infty) = O(n^{-\delta}), \ n \to \infty, \delta > 0. \]
This condition was considered in \cite[Proposition 5.7]{KN2018}. 
The simple random walk on $\mathbb{Z}^d, d \ge 3$, satisfies this condition. 

The following corresponds to \cite[Remark 3]{CCMP2026}. 
\begin{Prop}\label{prop:Lp-explode}
Assume that $(U_2)$ holds, that $0 < F_{\inf} \le F_{\sup} < 1$, and that $p \ge 1$. 
Let  $\alpha > p-2 + \frac{p-1}{\delta}$. 
Let $f(0) \coloneqq 0$ and $f(j) \coloneqq j^{\alpha/p} F_{\inf}^{- j/p}$ for $j \ge 1$. 
Then 
\[ \lim_{n \to \infty} E^x \left[ \left( \frac{G_n (f)}{n} \right)^p \right] = \infty. \]
\end{Prop}


\begin{proof}
We first remark that 
\[ E^x \left[ \left( \frac{G_n (f)}{n} \right)^p \right] \ge \frac{1}{n^p} \sum_{j \ge 1} f(j)^p E^x \left[\left(R_n^{(j)}\right)^p\right] \ge \frac{1}{n^p} \sum_{j \ge 2} f(j)^p E^x \left[R_n^{(j)}\right]. \]
By Lemma \ref{lem:exp-lower-quantitative}, 
\[ \frac{1}{n^p} \sum_{j \ge 2} f(j)^p E^x \left[R_n^{(j)}\right] \ge \frac{(1-F_{\sup})^2}{n^p}   \left\lfloor \frac{n}{2} \right\rfloor \sum_{j \ge 2}  j^{\alpha} F_{\inf}^{-j} F_{\inf}\left( \left\lfloor \frac{n}{2(j-1)} \right\rfloor\right)^{j-1}. \]
By  $(U_2)$, 
there exists $C > 0$ such that for every $n \ge 1$ and every $j \ge 2$, 
\[ F_{\inf}\left( \left\lfloor \frac{n}{2(j-1)} \right\rfloor\right) \ge  \left(1 - C \left(\frac{j-1}{n}\right)^{\delta}\right)_{+}  F_{\inf}. \]
Hence it suffices to show that 
\begin{equation}\label{eq:lim-infty-sum-j}
\lim_{n \to \infty} \frac{1}{n^{p-1}} \sum_{j \ge 2} j^{\alpha} \left(\left(1 - C \left(\frac{j-1}{n}\right)^{\delta}\right)_{+}\right)^{j-1} = \infty. 
\end{equation}

Let $a > 0$ such that $C a^{\delta + 1} < 1/2$. 
If $j - 1 \le \lfloor a n^{\delta/(1+\delta)} \rfloor$, then 
\[ \left(1 - C \left(\frac{j-1}{n}\right)^{\delta}\right)^{j-1}  \ge 1 - C \frac{(j-1)^{1+\delta}}{n^{\delta}} \ge \frac{1}{2}. \]
Hence
\[ \sum_{j \ge 2} j^{\alpha} \left(\left(1 - C \left(\frac{j-1}{n}\right)^{\delta}\right)_{+}\right)^{j-1} \ge \frac{1}{2} \sum_{j=2}^{\lfloor a n^{\delta/(1+\delta)} \rfloor + 1} j^{\alpha}  \]
and the right-hand side grows on the scale of $n^{(\alpha + 1) \delta / (\delta + 1)}$. 
By the assumption, $(\alpha + 1) \delta >  (p-1) (\delta + 1)$ and we obtain \eqref{eq:lim-infty-sum-j}. 
\end{proof}


\begin{Rem}
Suppose that the assumptions of Proposition \ref{prop:Lp-explode} hold. 
If, additionally, $p > 1$ and $F_{\sup} < F_{\inf}^{1/p}$, 
then $C_{\sup}(f) < \infty$ and 
\[ \lim_{n \to \infty} E^x \left[ \left( \frac{G_n (f)}{n} - C_{\sup}(f) \right)_{+} \right] =  \lim_{n \to \infty} E^x \left[ \left( \frac{G_n (f)}{n} - C_{\inf}(f) \right)_{-} \right] = 0. \] 
If $F_{\inf} = F_{\sup} < 1$ and $p >1$,  then $F_{\sup} < F_{\inf}^{1/p}$. 
\end{Rem}

In general, the strict inequality $F_{\inf} < F_{\sup}$ can occur. 
See \cite[Theorem 1.3]{Okamura2014}. 
By the same argument as in the proof of  \cite[Theorem 1.3]{Okamura2014}, 
one can show that 
\[ (F_{\inf})^{j-1} (1-F_{\sup}) =  \liminf_{n \to \infty} \frac{\inf_{x \in \mathcal{X}} E^{x} \left[R_n^j \right]}{n} \]
\[ < \limsup_{n \to \infty} \frac{\sup_{x \in \mathcal{X}} E^{x} \left[R_n^j \right]}{n} = (F_{\sup})^{j-1} (1-F_{\inf}) \]
for each $j$. 


\section{Examples}\label{sec:examples}


\begin{thebibliography}{99}
\bibitem[AK2020]{AK2020} Inna~M. Asymont and Dmitry Korshunov. \newblock Strong law of large numbers for a function of the local times of a transient random walk in {{\({\mathbb{Z}}^d\)}}. \newblock {\em J. Theor. Probab.}, 33(4):2315--2336, 2020.
\bibitem[CCMP2026]{CCMP2026} Yinshan Chang, Qinwei Chen, Qian Meng, and Xue Peng. \newblock Strong law of large numbers for a function of the local time of a transient random walk on a group. \newblock {\em J. Theor. Probab.}, 39(1):16, 2026. \newblock Id/No 7.
\bibitem[KN2018]{KN2018} Takashi Kumagai and Chikara Nakamura. \newblock Lamplighter random walks on fractals. \newblock {\em J. Theor. Probab.}, 31(1):68--92, 2018.
\bibitem[Okamura2014]{Okamura2014} Kazuki Okamura. \newblock On the range of random walk on graphs satisfying a uniform condition. \newblock {\em ALEA, Lat. Am. J. Probab. Math. Stat.}, 11(2):341--357, 2014.
\bibitem[Pitt1974]{Pitt1974} Joel~H. Pitt. \newblock Multiple points of transient random walks. \newblock {\em Proc. Am. Math. Soc.}, 43:195--199, 1974.
\bibitem[SCZ2016]{SCZ2016} Laurent Saloff-Coste and Tianyi Zheng. \newblock Random walks and isoperimetric profiles under moment conditions. \newblock {\em Ann. Probab.}, 44(6):4133--4183, 2016.
\bibitem[Woess2013]{Woess2013} Wolfgang Woess. \newblock What is a horocyclic product, and how is it related to lamplighters? \newblock {\em Int. Math. Nachr., Wien}, 224:1--27, 2013.
\end{thebibliography}
\end{document}
